%% kapa: exam=1 round=31 year=2020 period=1 form=A subject=civil_law subject_ko=민법 seq=9 number=9
%% source: Q-Net 감정평가사 기출문제 (cst00309) artlSeq=5207219 file=extracted/1008/5207219/A형/A형/2020년도 제31회 감정평가사 1차 시험_1교시_A형.pdf
%% conversion: kordoc 4.24.51 -> tools/gen1.py (auto); review=proofread-sonnet; flags=none
\begin{question}{9}
불공정한 법률행위에 관한 설명으로 옳은 것은? (다툼이 있으면 판례에 따름)
\begin{choices}
  \item 불공정한 법률행위로 무효가 된 행위의 전환은 인정되지 않는다.
  \item 불공정한 법률행위라도 당사자가 무효임을 알고 추인한 경우 유효로 될 수 있다.
  \item 불공정한 법률행위에 해당하는지 여부는 그 행위를 한 때를 기준으로 판단한다.
  \item 불공정한 법률행위의 요건을 갖추지 못한 법률행위는 반사회질서행위가 될 수 없다.
  \item 증여와 같이 아무런 대가관계 없이 당사자 일방이 상대방에게 일방적인 급부를 하는 행위도 불공정한 법률행위가 될 수 있다.
\end{choices}
\end{question}
\answer{3}
